\section{A C1 extension: $L$ a unit and $\eta$ infinitesimal}
\label{sec:r7-c1-L-unit-eta-small}

\begin{theorem}
The entire original C1 source subdomain
\[
 x=k^2\in\Inf_{>0},\qquad L=x\zeta\text{ a positive bounded unit},
 \qquad \eta\in\Inf_{>0}
\]
has a full-parent tensor-SOS certificate for every independent
projective testing pair.  All original openings and the retained
boundaries $D=0$ or $\beta=0$ remain allowed whenever they occur
in this source subdomain.
\end{theorem}

The proof is an exact adapter to the strong-II local and remote
certificates.  Its additional event is retained with a centered
square; it is not deleted from the parent.

\paragraph{The native source in the shared contact geometry.}
Keep the native notation of (NC1)--(NC4), but in this section denote
its $r,h,a$ by $r_n,h_n,a_n$ to avoid collisions.  Normalize $d=1$.
The exact native contact with $Z_c=1$ is
\[
 (u_c,v_c;W_c,X_c,Y_c,Z_c)
 =\left(\sqrt x,\sqrt{x\eta\zeta};
 \sqrt{\zeta/\eta},-\sqrt x\,h_n,
 \sqrt x\,r_n\sqrt{\zeta/\eta},1\right).
\]
It therefore has the same contact geometry used in (LII2), with
\[
 z=x,\quad y=L\eta,\quad B=L,\quad H=Lr_n=B-a_ny,
 \quad r=h_n=1-\theta,\quad t=a_n,\quad
 j=1+zr-t=J_n.                                           \tag{LC1}
\]
In particular $z,y$ are positive infinitesimals.  Here $r_n\sim1$,
$0<Lr_n<1$, and hence $b_0:=\operatorname{res}B$ satisfies
$0<b_0\le1$.  We use $b_0$ only for this source residue.

Set, as in the local strong-II proof,
\[
 c=Cyt,\quad d_p=Dyj,\quad
 \mu=m/z^2,\quad a_A=A/y,\quad
 \sigma=\alpha(1-H),\quad \tau=\beta(1-r)=\beta\theta.
\]
These are bounded source quantities, and $c,d_p,\mu,a_A,\sigma$
are positive units.  Writing $r_0=\operatorname{res}r$, exact
substitution in (NC2)--(NC4), with opening factors canceled first,
gives
\[
\begin{gathered}
 c,d_p\sim\frac{b_0}{2(1+b_0)},\qquad
 a_A\sim\frac1{2b_0(1+b_0)},\\
 \mu\sim\frac{r_0}{2b_0(1+b_0)},\qquad
 \sigma\sim\frac{b_0^2}{2(1+b_0)},\\
 \operatorname{res}\tau=
 \frac{b_0}{2(1+b_0)}-\frac{1-r_0}{2}.                    \tag{LC2}
\end{gathered}
\]
Since $\tau\ge0$, the last identity implies
$r_0\ge1/(1+b_0)>0$, proving the asserted unit lower bound for
$\mu$.  Also $t+j=1+zr\sim1$, and $c,d_p$ being units proves
$C,D$ unbounded uniformly through either opening $t=0$ or $j=0$
in residue.  Thus $D=0$ is in fact absent from this specific
subdomain by the exact source implications; it has not been excluded
by an additional hypothesis.  Literal $\beta=0$ is included.

The source-harmonic quantity $yCD/(C+D)$ and the cap behavior of
$\alpha$ have precisely the residues (RR1)--(RR2), with
$h=b_0$.  The remote proof (RR3)--(RR13) therefore applies unchanged,
retaining the extra positive packet
$4\beta D_v(X+uZ)^2$ throughout.  This covers every test outside
the local region where $u,v$ are both infinitesimal, including the
original projective endpoints.

\paragraph{The additional event and its exact local remainder.}
For the local region use exactly the graph, roots, harmonic pair,
and caps (LII2)--(LII4), with the source quantities (LC1)--(LC2).
Put
\[
 R=v(U-z)W,\qquad
 E_\beta=\xi+u\zeta+R/B
          =(X+uZ)-\frac{z(1-r)v}{B}W.                    \tag{LC3}
\]
This root vanishes at the actual contact.  Retain the entire
positive packet $4\beta D_vE_\beta^2$, in addition to the first,
second and fourth packets of (LII5); there is no $n$ packet since
$n=0$.  The exact remainder from the $\beta$ event is
\[
 \begin{split}
 \Delta_\beta={}&
 \frac{8\tau z}{B}D_vvW(\xi+u\zeta)\\
 &+\frac{4\tau z}{B^2}D_vV\{2U-z(1+r)\}W^2.
 \end{split}                                             \tag{LC4}
\]
All its coefficients are bounded even if $\beta$ is unbounded:
only $\tau=\beta(1-r)$ occurs.  The exact full remainder is (LII6)
with its $n$ line deleted, plus (LC4).  Thus no state normal,
boundary coefficient, or first contact jet has been omitted.

The $R_hW$ gauge (LII7) is unchanged, because (LC4) contains no
$R_h$.  The coefficient representation (LII8), its exact contact
and first-normal vanishings (LII9)--(LII10), and its bounded monic
divisions (LII12)--(LII13) all remain valid.  Explicitly, (LC4)
adds $8\tau z(1+V)/B$ to $P_\xi$ and $P_\zeta$, adds zero to
$P_N$, and adds its printed scalar line to $\phi$.  Hence the
essential support $P_\xi\in(U,z)$ and $\deg_U\phi\le2$ remain
literal polynomial statements.

The only changed data in (LII11) are
\[
\begin{split}
 d_s&=\frac{4z^2}{B^2}
   \left\{\mu(B+zr)^2+\frac{r^2}{2}-\tau(1+r)\right\},\\
 e_s&=\frac{4z^2}{B^2}
   \left\{\frac{r^2}{2}-\tau(1+r)\right\}.
\end{split}                                             \tag{LC5}
\]
Both quotients by $z^2$ are bounded; $e_s$ need not be positive.
Its contribution to the final Gram has an infinitesimal coefficient
because it is multiplied by $V$.  The decisive exact residue
cancellation, using (LC2), is
\[
 \begin{split}
 \operatorname{res}\left\{
 \mu B^2+\frac{r^2}{2}-\tau(1+r)\right\}
 &=\frac{b_0r_0}{2(1+b_0)}+\frac{r_0^2}{2}
 -\left\{\frac{b_0}{2(1+b_0)}-\frac{1-r_0}{2}\right\}(1+r_0)\\
 &=\frac1{2(1+b_0)}.
 \end{split}                                             \tag{LC6}
\]
Consequently
\[
 \operatorname{res}(d_s/z^2)=\frac2{b_0^2(1+b_0)},\qquad
 \operatorname{res}A_\xi=\frac4{b_0(1+b_0)}.              \tag{LC7}
\]
The second equality is unchanged from (LII18): the new $P_\xi$
term is a bounded multiple of $z$ independent of $U$, so it has
zero $U$-quotient residue.  The four hard-normal residues (LII17)
are unchanged as well.

Apply the exact allocations (LII15) to all the coefficients,
including (LC4).  They give the same bounded six-root Gram
$(N,\xi,\zeta,R_h,R,T)$ as in (LII16).  Its finite $(\xi,R)$
block is exactly (LII19), with $h=b_0$, and has strictly positive
Schur complement $2/[b_0(1+b_0)^2]$.  The same exact conductor
(LII20) supplies an arbitrarily large real $T^2$ pivot while
changing every other residue entry by zero.  The complete bounded
Gram is strict.  Restoring the four retained packets, including
$4\beta D_vE_\beta^2$, proves the local tensor SOS for all
independent rates and literal affine axes.  The remote admission
above completes the theorem.

\paragraph{Verification and exact frontier effect.}
The accompanying checker reconstructs the entire original C1 parent
from this remainder and all four retained packets, verifies the new
centered-event identity, the unchanged contact-ideal support, every
new source residue, the finite block, and the conductor.  Stationarity
itself is the unchanged native source identity, not an added
restriction.  The result removes the whole $L$-unit, $\eta$-infinitesimal
subdomain from C1.4.  It does not assert the other $L$-unit source
regimes or the entire C1.4 row.
